% Mkusanyiko wa mazoezi — somo la 2
% Somo hutambuliwa kutokana na jina la faili hii.

\begin{exo}[Kutoka barafu baridi hadi maji yanayochemka: makadirio ya kiwango]{chauffage-glace-eau-ebullition}
\seoready{true}
\keywords{kalorimetria, mabadiliko ya hali, kuyeyuka, uvukizaji, joto fiche, kiwango cha ukubwa}

\textbf{Data za namba (kwenye shinikizo la angahewa):}
\[
\begin{aligned}
c_{\text{barafu}} &= 2{,}1\times10^3\,\text{J}\!\cdot\!\text{kg}^{-1}\!\cdot\!\text{K}^{-1},
& L_f &= 3{,}34\times10^5\,\text{J}\!\cdot\!\text{kg}^{-1},\\
c_{\text{maji}} &= 4{,}18\times10^3\,\text{J}\!\cdot\!\text{kg}^{-1}\!\cdot\!\text{K}^{-1},
& L_v &= 2{,}26\times10^6\,\text{J}\!\cdot\!\text{kg}^{-1}.
\end{aligned}
\]
Kwenye shinikizo la angahewa, kilogramu moja ya barafu yenye halijoto ya mwanzo $-100\,{}^\circ\text{C}$ hupashwa hatua kwa hatua.

\begin{enumerate}
  \item Kokotoa nishati inayohitajika kuleta barafu kutoka $-100\,{}^\circ\text{C}$ hadi $0\,{}^\circ\text{C}$.
  \item Kokotoa nishati inayohitajika kuyeyusha barafu yote kwenye $0\,{}^\circ\text{C}$.
  \item Kokotoa nishati inayohitajika kupasha maji ya kioevu kutoka $0\,{}^\circ\text{C}$ hadi $100\,{}^\circ\text{C}$.
  \item Kokotoa nishati ya ziada inayohitajika kuvukiza maji yote kwenye $100\,{}^\circ\text{C}$.
  \item Ni hatua ipi inayohitaji nishati nyingi zaidi? Linganisha nishati ya kupasha kilogramu moja ya maji kutoka $0\,{}^\circ\text{C}$ hadi $100\,{}^\circ\text{C}$ na nishati tarajiwa ya uvutano ya misa $m$ inayoanguka kutoka urefu wa mita kumi. Ni misa gani $m$ inapaswa kuangushwa ili kutoa nishati hiyo?
  \item Sasa zingatia mchakato wa kinyume. Misa $m_v$ ya mvuke wa maji huganda kuwa kimiminika kwenye kioo cha mbele cha gari, bila maji yaliyopatikana kuendelea kupoa. Andika joto $Q_{\text{tolewa}}$ linalotolewa wakati wa ugandishaji huo, kisha ulikokotoe kwa $m_v=1{,}0\times10^{-2}\,\text{kg}$. Kwenye halijoto ya kioo tumia $L_v=2{,}45\times10^6\,\text{J}\!\cdot\!\text{kg}^{-1}$.
\end{enumerate}

\begin{indication}
Kwa kupasha hali moja bila mabadiliko ya hali, tumia $Q=mc\Delta T$. Kwa mabadiliko ya hali kwenye halijoto isiyobadilika, tumia $Q=mL$. Nishati tarajiwa ya uvutano ya misa $m$ kwenye urefu $h$ ni $E_p=mgh$; tumia $g=9{,}81\,\text{m}\!\cdot\!\text{s}^{-2}$.
\end{indication}

\begin{solution}
\textbf{1.} Kupasha barafu hadi kiwango chake cha kuyeyuka kunahitaji
\[
Q_1=mc_{\text{barafu}}\Delta T
=1{,}0\times2{,}1\times10^3\times100
=2{,}1\times10^5\,\text{J}.
\]

\textbf{2.} Wakati wa kuyeyuka halijoto hubaki $0\,{}^\circ\text{C}$:
\[
Q_2=mL_f=1{,}0\times3{,}34\times10^5=3{,}34\times10^5\,\text{J}.
\]

\textbf{3.} Kupasha maji ya kioevu kutoka $0$ hadi $100\,{}^\circ\text{C}$ kunahitaji
\[
Q_3=mc_{\text{maji}}\Delta T
=1{,}0\times4{,}18\times10^3\times100
=4{,}18\times10^5\,\text{J}.
\]

\textbf{4.} Uvukizaji kamili unahitaji tena
\[
Q_4=mL_v=1{,}0\times2{,}26\times10^6=2{,}26\times10^6\,\text{J}.
\]
Hatua hii pekee inahitaji nishati ya kiwango cha megajouli kadhaa.

\textbf{5.} Uvukizaji ndio hatua inayohitaji nishati nyingi zaidi:
\[
Q_4=2{,}26\times10^6\,\text{J}
>Q_3=4{,}18\times10^5\,\text{J}
>Q_2=3{,}34\times10^5\,\text{J}
>Q_1=2{,}1\times10^5\,\text{J}.
\]
Uvukizaji huhitaji takribani $(2{,}26\times10^6)/(4{,}18\times10^5)\approx5{,}4$ zaidi ya kupasha maji ya kioevu kutoka $0$ hadi $100\,{}^\circ\text{C}$.

Kwa misa $m$ inayoanguka kutoka $h=10\,\text{m}$, $E_p=mgh$. Kwa kuweka $mgh=Q_3$ tunapata
\[
m=\frac{Q_3}{gh}
=\frac{4{,}18\times10^5}{9{,}81\times10}
\approx4{,}26\times10^3\,\text{kg}.
\]
Kwa hiyo ingelazimu kuangusha karibu tani $4{,}3$ kutoka mita kumi ili kutoa nishati sawa na ile ya kupasha kilogramu moja ya maji kutoka $0$ hadi $100\,{}^\circ\text{C}$! Hii inaeleza kwa nini kupasha maji ni sehemu muhimu ya matumizi ya nishati ya nyumbani. Uwezo mahususi wa joto wa maji pia ni mkubwa sana ukilinganishwa na metali za kawaida; kwa mfano ni karibu mara kumi ule wa shaba (tazama mazoezi yanayofuata).

\textbf{6.} Ugandishaji wa mvuke kuwa kioevu ni kinyume cha uvukizaji: mvuke hutoa joto fiche kwa kioo,
\[
Q_{\text{tolewa}}=m_vL_v.
\]
Kwa $m_v=1{,}0\times10^{-2}\,\text{kg}$,
\[
Q_{\text{tolewa}}=1{,}0\times10^{-2}\times2{,}45\times10^6
=2{,}45\times10^4\,\text{J}.
\]
Hivyo hata kiasi kidogo cha mvuke kinaweza kutoa joto lisilopuuzika kinapoganda kuwa kimiminika. Joto hili hupokelewa na kioo na mazingira yake.
\end{solution}
\end{exo}

\begin{exo}[Uvukizaji na utoaji wa mvuke wa mti mkubwa: kiyoyozi asilia?]{odg-evapotranspiration-arbre}
\seoready{true}
\keywords{kiwango cha ukubwa, uvukizaji na utoaji wa mvuke, mti, joto fiche, nguvu ya upoozaji, kiyoyozi, COP}

Katika siku yenye joto na ukavu, mti mkubwa wenye maji ya kutosha unaweza kutoa kwa uvukizaji na utoaji wa mvuke kiasi cha karibu $500\,\text{L}=5{,}0\times10^{-1}\,\text{m}^3$. Kwa kuwa utoaji wa mvuke hutokea hasa kwenye mwanga, tutadhani unasambaa kwa saa kumi na mbili. Majani ya mti yanafunika eneo la ardhini $A_{\text{mti}}=50\,\text{m}^2$. Kwenye shinikizo la angahewa, densiti ya maji na joto fiche la uvukizaji ni
\[
\rho_{\text{maji}}=1{,}0\times10^3\,\text{kg}\!\cdot\!\text{m}^{-3},
\qquad L_v=2{,}45\times10^6\,\text{J}\!\cdot\!\text{kg}^{-1}.
\]

\begin{enumerate}
  \item Kokotoa misa ya maji inayotolewa na mti kwa siku, kisha nishati ya kuyavukiza.
  \item Eleza kwa nini jambo hili linaleta athari ya upoozaji.
  \item Kokotoa nguvu wastani ya upoozaji ya mti kwa kila mita ya mraba katika saa kumi na mbili za utoaji hai wa mvuke.
  \item Kwa kulinganisha, kiyoyozi kinatumia nguvu ya umeme $P_{\mathrm{elec}}=2{,}0\times10^3\,\text{W}$, kina mgawo wa utendaji $\mathrm{COP}=3{,}0$, na kupooza chumba cha eneo $A_{\text{chumba}}=20\,\text{m}^2$. Kwa ufafanuzi,
\[
\mathrm{COP}=\frac{P_{\text{baridi}}}{P_{\mathrm{elec}}},
\]
ambapo $P_{\text{baridi}}$ ni nguvu ya joto inayoondolewa chumbani. Kokotoa $P_{\text{baridi}}$ na nguvu yake kwa kila mita ya mraba, kisha linganisha na mti.
  \item Kwa nini hatuwezi kuhitimisha kuwa mti na kiyoyozi hutoa upoozaji unaohisiwa sawa kabisa hata kama nguvu zao kwa kila eneo zina kiwango kimoja cha ukubwa?
  \item Je, kwa maoni yako, mimea mingi ya ndani inaweza kupooza nyumba? (Maarifa ya jumla: jibu hutegemea michakato ya kibiolojia.)
\end{enumerate}

\begin{indication}
Tumia $m=\rho V$, $Q_{\mathrm{evap}}=mL_v$ na $P=Q/\tau$. Kwa kiyoyozi, $P_{\text{baridi}}=\mathrm{COP}\,P_{\mathrm{elec}}$.
\end{indication}

\begin{solution}
\textbf{1.} $500\,\text{L}=5{,}0\times10^{-1}\,\text{m}^3$, kwa hiyo
\[
m=\rho_{\text{maji}}V=1{,}0\times10^3\times5{,}0\times10^{-1}=5{,}0\times10^2\,\text{kg},
\]
na
\[
Q_{\mathrm{evap}}=mL_v=5{,}0\times10^2\times2{,}45\times10^6=1{,}225\times10^9\,\text{J}.
\]
Kiwango chake ni gigajouli moja.

\textbf{2.} Uvukizaji ni mabadiliko ya hali yanayofyonza joto. Nishati $Q_{\mathrm{evap}}$ hutolewa na majani, udongo na hewa inayozunguka; kupoteza joto huku hupunguza halijoto yao.

\textbf{3.} Saa kumi na mbili ni $\tau=12\times3600=4{,}32\times10^4\,\text{s}$. Hivyo
\[
P_{\text{mti}}=\frac{Q_{\mathrm{evap}}}{\tau}
=\frac{1{,}225\times10^9}{4{,}32\times10^4}
\approx2{,}84\times10^4\,\text{W},
\]
na kwa $50\,\text{m}^2$,
\[
\frac{P_{\text{mti}}}{A_{\text{mti}}}\approx5{,}7\times10^2\,\text{W}\!\cdot\!\text{m}^{-2}.
\]

\textbf{4.}
\[
P_{\text{baridi}}=\mathrm{COP}\,P_{\mathrm{elec}}=3{,}0\times2{,}0\times10^3=6{,}0\times10^3\,\text{W},
\]
na
\[
\frac{P_{\text{baridi}}}{A_{\text{chumba}}}=\frac{6{,}0\times10^3}{20}=3{,}0\times10^2\,\text{W}\!\cdot\!\text{m}^{-2}.
\]
Katika modeli hii, uvukizaji na utoaji wa mvuke wa mti una nguvu karibu mara mbili ya kiyoyozi cha kawaida.

\textbf{5.} Ulinganisho unahusu tu mtiririko wa nishati. Kiyoyozi hupooza ujazo uliofungwa na kudhibitiwa, lakini mvuke wa mti hutawanyika angani na upepo huleta hewa moto. Athari inayohisiwa hutegemea uingizaji hewa, unyevunyevu, halijoto ya mazingira na umbali kutoka kwa mti. Utoaji wa mvuke pia hubadilika kwa mwanga, upepo, unyevunyevu, spishi na maji yaliyopo, na unaweza kupungua sana wakati wa ukame. Kwa upande mwingine, mti pia hutoa kivuli. Kwa hiyo hesabu hizi hulinganisha kiwango cha ukubwa, si usawa kamili wa faraja ya joto.

\textbf{6.} Kwa vitendo, upoozaji wa mimea ya ndani ungekuwa mdogo sana. Kwa mimea mingi, mwanga hufungua stomata ambazo mvuke wa maji hupitia; mwanga wa ndani kwa kawaida ni hafifu, hivyo stomata hufunguka kidogo. Katika chumba kisicho na uingizaji hewa mzuri, unyevunyevu unaoongezeka pia hupunguza uvukizaji. Mimea mingi inaweza kutoa upoozaji mdogo wa karibu, lakini haiwezi kuchukua nafasi ya kiyoyozi. Hewa yenye unyevunyevu ingepaswa kupelekwa nje. Mwanga mkali wa bandia unaweza kuongeza utoaji wa mvuke, lakini nishati yake ya umeme hatimaye hubadilika karibu yote kuwa joto chumbani. Mimea ya CAM ya mazingira kavu, kama kaktasi, ni tofauti: stomata zake hufunguka hasa usiku.
\end{solution}
\end{exo}

\begin{exo}[Kuchanganya misa mbili za maji]{calorimetrie-melange-eau}
\seoready{true}
\keywords{kalorimetria, mchanganyiko, uwezo wa joto, uwezo mahususi wa joto, Joseph Black, halijoto ya usawa}

Misa $m_1=0{,}200\,\text{kg}$ ya maji kwenye $T_1=80\,{}^\circ\text{C}$ humiminwa kwenye chombo chenye tayari misa $m_2=0{,}300\,\text{kg}$ ya maji kwenye $T_2=15\,{}^\circ\text{C}$. Chombo ni kalorimita bora: tunapuuza kupotea kwa joto kwenda nje na uwezo wa joto wa chombo chenyewe. Kumbuka
\[
Q=mc\Delta T,
\]
ambapo uwezo mahususi wa joto wa maji ya kioevu ni $c=4{,}18\times10^3\,\text{J}\!\cdot\!\text{kg}^{-1}\!\cdot\!\text{K}^{-1}$.

\begin{enumerate}
  \item Thibitisha kuwa joto lote linalotolewa na maji moto hupokelewa na maji baridi; andika mlinganyo wa uhifadhi wenye $m_1$, $m_2$, $c$, $T_1$, $T_2$ na halijoto ya usawa $T_{eq}$.
  \item Pata fomula ya $T_{eq}$, kisha thamani yake ya namba.
  \item Kokotoa joto $Q$ linalotolewa na maji moto wakati wa kuchanganya.
  \item Je, mchanganyiko kama huu wa dutu ileile unaweza kutumika kuamua $c$? Eleza.
\end{enumerate}

\begin{indication}
Kwa kuwa kalorimita imetengwa na ni ngumu, nishati ya ndani ya jumla $U_1+U_2$ huhifadhiwa: joto linalotolewa na moja hupokelewa na nyingine kwa ishara kinyume.
\end{indication}

\begin{solution}
\textbf{1.} Katika kalorimita iliyotengwa, nishati ya ndani ya mfumo $\{$maji moto, maji baridi$\}$ huhifadhiwa. Kwa kutumia $Q=mc\Delta T$ kwa kila misa kutoka hali ya mwanzo hadi usawa,
\[
m_1c(T_{eq}-T_1)+m_2c(T_{eq}-T_2)=0.
\]

\textbf{2.} Kwa kuwa ni dutu ileile, $c$ hufutika:
\[
T_{eq}=\frac{m_1T_1+m_2T_2}{m_1+m_2}
=\frac{0{,}200\times80+0{,}300\times15}{0{,}200+0{,}300}
=41\,{}^\circ\text{C}.
\]

\textbf{3.}
\[
Q=m_1c(T_1-T_{eq})
=0{,}200\times4{,}18\times10^3\times(80-41)
\approx3{,}26\times10^4\,\text{J}.
\]
Maji baridi hupokea kiasi kilekile:
\[
m_2c(T_{eq}-T_2)=0{,}300\times4{,}18\times10^3\times26\approx3{,}26\times10^4\,\text{J}.
\]

\textbf{4.} La. Kwa misa mbili za dutu ileile, $c$ ileile huzidisha vipengele vyote vya mizania na hufutika. Halijoto ya usawa haitegemei $c$; ni wastani wa halijoto za mwanzo uliopimwa kwa misa. Ili kupima uwezo mahususi wa joto kwa mbinu ya mchanganyiko, lazima tuongeze dutu yenye uwezo wa joto unaojulikana, kama katika zoezi linalofuata.
\end{solution}
\end{exo}

\begin{exo}[Kuamua uwezo mahususi wa joto wa metali kwa mbinu ya mchanganyiko]{calorimetrie-methode-melanges}
\seoready{true}
\keywords{kalorimetria, mbinu ya mchanganyiko, uwezo mahususi wa joto, kalorimita, joto mahususi, thamani sawia ya maji}

Tunataka kupima uwezo mahususi wa joto $c_m$ wa metali isiyojulikana kwa \emph{mbinu ya mchanganyiko}, kama Joseph Black alivyofanya katika karne ya XVIII\textsuperscript{e}. Sampuli yenye misa $m=0{,}150\,\text{kg}$ hupashwa hadi $T_m=95\,{}^\circ\text{C}$, kisha huzamishwa haraka katika kalorimita yenye $m_e=0{,}250\,\text{kg}$ za maji kwenye $T_e=18\,{}^\circ\text{C}$; $c_e=4{,}18\times10^3\,\text{J}\!\cdot\!\text{kg}^{-1}\!\cdot\!\text{K}^{-1}$. Halijoto ya usawa ni $T_{eq}=22\,{}^\circ\text{C}$. Mwanzoni puuza uwezo wa joto wa chombo (kuta, kichanganyio na kipimajoto).

\begin{enumerate}
  \item Kama katika zoezi lililotangulia, andika uhifadhi wa nishati wa mfumo uliotengwa $\{$metali, maji$\}$, ukiwa na $c_m$ kama kisichojulikana pekee.
  \item Pata fomula na thamani ya namba ya $c_m$. Inalingana takribani na metali ipi ya kawaida? Thamani za marejeo: alumini $\approx9{,}0\times10^2\,\text{J}\!\cdot\!\text{kg}^{-1}\!\cdot\!\text{K}^{-1}$, chuma $\approx4{,}5\times10^2$, shaba $\approx3{,}9\times10^2$, risasi $\approx1{,}3\times10^2$.
  \item Kwa kweli chombo chenyewe hufyonza sehemu ya joto la metali. Tunawakilisha athari hii kwa \emph{thamani sawia ya maji} $\mu=1{,}5\times10^{-2}\,\text{kg}$: kalorimita hufanya kana kwamba ni misa ya ziada $\mu$ ya maji. Kokotoa tena $c_m$ ukijumuisha $\mu$, na eleza mwelekeo wa sahihisho (je, $c_m$ iliyosahihishwa ni kubwa au ndogo kuliko ya swali la 2?).
  \item Kwa nini ni muhimu kuzamisha sampuli \emph{haraka} na kuitenga vizuri na hewa ya mazingira wakati wa kipimo?
\end{enumerate}

\begin{indication}
Metali hutoa joto $Q_m=mc_m(T_m-T_{eq})$, linalopokelewa lote na maji (na katika swali la 3, kalorimita inayolinganishwa na misa ya ziada $\mu$ ya maji): $Q_m=(m_e+\mu)c_e(T_{eq}-T_e)$.
\end{indication}

\begin{solution}
\textbf{1.} Katika mfumo uliotengwa, joto linalotolewa na metali ikipoa kutoka $T_m$ hadi $T_{eq}$ hupokelewa lote na maji yanayopashwa kutoka $T_e$ hadi $T_{eq}$:
\[
mc_m(T_m-T_{eq})=m_ec_e(T_{eq}-T_e).
\]

\textbf{2.}
\[
c_m=\frac{m_ec_e(T_{eq}-T_e)}{m(T_m-T_{eq})}
=\frac{0{,}250\times4{,}18\times10^3\times(22-18)}{0{,}150\times(95-22)}
\approx3{,}82\times10^2\,\text{J}\!\cdot\!\text{kg}^{-1}\!\cdot\!\text{K}^{-1}.
\]
Thamani hii iko karibu sana na ya shaba ($\approx3{,}9\times10^2\,\text{J}\!\cdot\!\text{kg}^{-1}\!\cdot\!\text{K}^{-1}$).

\textbf{3.} Maji na chombo sasa hupokea joto pamoja:
\[
c_m=\frac{(m_e+\mu)c_e(T_{eq}-T_e)}{m(T_m-T_{eq})}
=\frac{(0{,}250+0{,}015)\times4{,}18\times10^3\times4}{0{,}150\times73}
\approx4{,}05\times10^2\,\text{J}\!\cdot\!\text{kg}^{-1}\!\cdot\!\text{K}^{-1}.
\]
Sahihisho huongeza $c_m$ kidogo. Kupuuza thamani sawia ya maji katika swali la 2 kulipunguza makadirio ya joto halisi lililotolewa na metali—sehemu yake ilipasha chombo—na hivyo kulipunguza makadirio ya $c_m$.

\textbf{4.} Mbinu inategemea mfumo kubaki umetengwa wakati wa kipimo. Kuzamisha sampuli haraka hupunguza muda wa kubadilishana joto na hewa kabla haijafikia maji; kutenga kalorimita hupunguza upotevu nje wakati usawa unaanzishwa. Uhamisho wowote wa joto usiohesabiwa ungeharibu mizania ya swali la 1 na $c_m$ inayotokana nayo. Huo ndio uangalifu wa majaribio ambao Black, na baadaye Lavoisier na Laplace, walitumia katika kalorimita yao ya barafu.
\end{solution}
\end{exo}

\begin{exo}[Kalori za chakula na nguvu ya kimetaboliki]{calorimetrie-calories-alimentaires}
\seoready{true}
\keywords{kalori, kilokalori, Kalori ya chakula, kilinganishi cha kimekanika, nguvu ya kimetaboliki, vipimo}

Lebo ya lishe inaonyesha kuwa mtu mzima hutumia wastani wa karibu $2000\,\text{kcal}$ kwa siku. Kilokalori ni kipimo kisicho cha SI ambacho bado hutumika katika lishe; ubadilishaji wake ni $1\,\text{kcal}=4{,}18\times10^3\,\text{J}$.

\begin{enumerate}
  \item Badilisha kiasi hiki cha nishati ya kila siku kuwa jouli.
  \item Pata nguvu wastani ya kimetaboliki ya mtu huyu kwa wati, ukidhani nishati inatumika sawasawa katika saa 24.
  \item Baa ya nishati imeandikwa ``$210\,\text{kcal}$''. Ikiwa nishati hiyo ingebadilishwa yote kuwa joto, ni misa gani ya maji kwenye $20\,{}^\circ\text{C}$ ingeweza kufikishwa kwenye kiwango cha kuchemka ($100\,{}^\circ\text{C}$)? Tumia $c_{\text{maji}}=4{,}18\times10^3\,\text{J}\!\cdot\!\text{kg}^{-1}\!\cdot\!\text{K}^{-1}$.
  \item Kwa maoni yako, nishati hii hutumiwa kwa maumbo gani hasa katika mwili wa binadamu? (jibu la kimaelezo)
\end{enumerate}

\begin{indication}
Kwa swali la 2, tumia $\mathcal P=E/\Delta t$. Kwa swali la 3, kwanza badilisha nishati ya lebo kuwa jouli, kisha tenga $m$ katika $Q=mc\Delta T$.
\end{indication}

\begin{solution}
\textbf{1.} Katika vipimo vya SI,
\[
E=2000\times4{,}18\times10^3=8{,}36\times10^6\,\text{J}.
\]

\textbf{2.}
\[
\mathcal P=\frac{8{,}36\times10^6}{8{,}64\times10^4}\approx97\,\text{W}.
\]
Kwa siku nzima, binadamu hutumia kwa wastani karibu nguvu sawa na balbu ya kawaida ya nyuzi. Kiwango hiki kinachoshangaza kinaonyesha umuhimu wa kubadilisha kalori za chakula kuwa vipimo vya kawaida vya fizikia.

\textbf{3.} Nishati ya baa ni
\[
Q=210\times4{,}18\times10^3\approx8{,}78\times10^5\,\text{J}.
\]
Kwa $\Delta T=80\,\text{K}$,
\[
m=\frac{Q}{c\Delta T}
=\frac{8{,}78\times10^5}{4{,}18\times10^3\times80}
\approx2{,}63\,\text{kg}.
\]
Kwa hiyo, kwa kiwango cha ukubwa, baa moja ina nishati ya kuleta karibu $2{,}5\,\text{kg}$ za maji ya bomba kwenye kiwango cha kuchemka.

\textbf{4.} Mwili hauondoi nishati inaposemwa ``unaitumia''; hubadilisha nishati ya kemikali ya virutubisho. Sehemu hufanya kazi ya kimekanika ya misuli, kuendesha viungo, kudumisha tofauti za umeme za seli, kusafirisha molekuli na kufanya usanisi wa kemikali wa kurekebisha tishu. Sehemu nyingine inaweza kuhifadhiwa kikemikali katika glikojeni na mafuta. Hatimaye sehemu kubwa ya nishati iliyotumiwa hutawanywa kama joto, ambalo husaidia kudumisha halijoto ya mwili kabla ya kuhamishwa mazingira. Sehemu ndogo pia hutoka kama nishati ya kemikali iliyobaki katika taka.
\end{solution}
\end{exo}
